\paragraph{Jensen 递归阈值、primitive twist 资源律与 prime-slice 支撑障碍}
在 Jensen 缺陷的帽核二阶差分、圆盘完成化的 RH 证书、单位根扭曲的四阶慢模展开以及 prime-slice 外置与局部化 solenoid 对偶都已建立之后，还可把这些接口进一步压缩为一条共同的有限层刚性链。其核心现象分别是：中心单零点在 dyadic 递归中具有 sharp 的噪声阈值；主 primitive twist 的抑制代价必服从二次资源律；而一切处处分支的前史链一旦要求忠实 prime-slice 外置，便不可能压入固定有限支持的局部化 solenoid。

\begin{theorem}[中心单零点的 dyadic Jensen 递归阈值与自证化]\label{thm:xi-time-part9zq-jensen-central-root-threshold}
\leanverified{paper\_xi\_time\_part9zq\_jensen\_central\_root\_threshold}
在定理~\ref{thm:xi-jensen-defect-hat-kernel-second-difference} 与定理~\ref{thm:xi-jensen-recursive-resolution-noise-coupling} 的设定下，设 \(x_\ast\) 是 \(J_F\) 对应的一个中心单零点位置，并对某个尺度 \(h>0\) 给出带噪三点观测
\[
J^\varepsilon(x_\ast-h),\qquad
J^\varepsilon(x_\ast),\qquad
J^\varepsilon(x_\ast+h),
\]
满足
\[
\left|J^\varepsilon(\xi)-J_F(\xi)\right|\le \varepsilon
\qquad
\bigl(\xi\in\{x_\ast-h,x_\ast,x_\ast+h\}\bigr).
\]
记
\[
W_h^\varepsilon(x_\ast)
:=
J^\varepsilon(x_\ast+h)-2J^\varepsilon(x_\ast)+J^\varepsilon(x_\ast-h).
\]
则成立：
\begin{enumerate}
  \item 若
  \[
  \varepsilon<\frac{h}{4},
  \]
  则
  \[
  W_h^\varepsilon(x_\ast)>0;
  \]
  因而以 \(W_h^\varepsilon>0\) 为触发规则的中心化 dyadic 递归在该层必然继续；
  \item 若更强地
  \[
  \varepsilon<\frac{h}{8},
  \]
  则
  \[
  W_h^\varepsilon(x_\ast)>4\varepsilon,
  \]
  因而该层已构成一个有限三点零点证书，足以推出窗口 \((x_\ast-h,x_\ast+h)\) 内确有盘内零点；
  \item 常数 \(1/4\) 对“继续触发”这一目的而言是尖锐的。
\end{enumerate}
\end{theorem}
\begin{proof}
由定理~\ref{thm:xi-jensen-recursive-resolution-noise-coupling}，在中心单零点情形下，
\[
\varepsilon<\frac{h}{4}
\quad\Longrightarrow\quad
W_h^\varepsilon(x_\ast)>0,
\]
且该门槛对触发规则 \(W_h^\varepsilon>0\) 是尖锐的。这给出第 \(1\) 条与第 \(3\) 条。

若 \(\varepsilon<h/8\)，则
\[
W_h^\varepsilon(x_\ast)\ge h-4\varepsilon>4\varepsilon.
\]
于是可应用推论~\ref{cor:xi-jensen-defect-noisy-threepoint-certificate}，得到
\[
\mu_F((x_\ast-h,x_\ast+h))\ge 1,
\]
即对应半径带内必有盘内零点。证毕。
\end{proof}

\begin{corollary}[无穷深度 dyadic 细化的临界噪声律]\label{cor:xi-time-part9zq-jensen-dyadic-critical-noise}
沿用定理~\ref{thm:xi-time-part9zq-jensen-central-root-threshold} 的设定。令
\[
h_n:=2^{-n}h_0,
\qquad
\varepsilon_n\ge 0
\qquad(n\ge 0).
\]
则：
\begin{enumerate}
  \item 若
  \[
  \sup_{n\ge 0}\frac{4\varepsilon_n}{h_n}<1,
  \]
  则中心化递归在全部尺度上永不熄灭；
  \item 若更强地
  \[
  \sup_{n\ge 0}\frac{8\varepsilon_n}{h_n}<1,
  \]
  则全部尺度同时自证化：每一层都给出一个有限三点零点证书；
  \item 反之，只要存在某层 \(n_0\) 使
  \[
  4\varepsilon_{n_0}\ge h_{n_0},
  \]
  则不能对一切合法噪声实现统一保证该层继续触发。
\end{enumerate}
\end{corollary}
\begin{proof}
前两条逐层应用定理~\ref{thm:xi-time-part9zq-jensen-central-root-threshold} 即得。第三条则是同一定理中 sharpness 结论在第 \(n_0\) 层的直接重述。证毕。
\end{proof}

\begin{theorem}[单个三点 Jensen 见证给出的有限 RH 否证器]\label{thm:xi-time-part9zq-threepoint-jensen-rh-disproof}
\leanverified{paper\_xi\_time\_part9zq\_threepoint\_jensen\_rh\_disproof}
令
\[
F_\zeta(w):=\Xi_\zeta(s(w))
\]
为圆盘完成化读出函数，并记其 Jensen 缺陷为 \(J_\zeta\)。若对某个中心 \(x\)、半宽 \(h\) 与噪声水平 \(\varepsilon\) 已知
\[
\left|J_\zeta^\varepsilon(\xi)-J_\zeta(\xi)\right|\le \varepsilon
\qquad
\bigl(\xi\in\{x-h,x,x+h\}\bigr),
\]
且
\[
W_h^\varepsilon(x)
:=
J_\zeta^\varepsilon(x+h)-2J_\zeta^\varepsilon(x)+J_\zeta^\varepsilon(x-h)
>
4\varepsilon,
\]
则对应半径窗口内必存在 \(F_\zeta\) 的盘内零点。等价地，视界对数导数
\[
\mathscr C_\zeta(w):=-\frac{\Xi_\zeta'(s(w))}{\Xi_\zeta(s(w))}
\]
在圆盘内部存在极点。因而 RH 失效。
\end{theorem}
\begin{proof}
由推论~\ref{cor:xi-jensen-defect-noisy-threepoint-certificate}，条件
\[
W_h^\varepsilon(x)>4\varepsilon
\]
直接推出该半径窗口对应的对数半径带中至少存在一个盘内零点，即 \(F_\zeta\) 在 \(\mathbb D\) 内有零点。

再由定理~\ref{thm:xi-disk-zero-pole-index-reduction}，
\[
F_\zeta(w_0)=0
\iff
\mathscr C_\zeta \text{ 在 }w_0\text{ 处有极点},
\]
而附录定理~\ref{thm:app-rh-iff-disk-zero-free} 又给出 RH 等价于 \(F_\zeta\) 在单位圆盘内无零点。故一旦存在上述三点见证，便得到 RH 失效。证毕。
\end{proof}

\begin{corollary}[单个三点 Jensen 见证强制负平方指数为正]\label{cor:xi-time-part9zq-threepoint-jensen-negative-square-index}
在定理~\ref{thm:xi-time-part9zq-threepoint-jensen-rh-disproof} 的设定下，一旦存在某个三点噪声见证满足
\[
W_h^\varepsilon(x)>4\varepsilon,
\]
则
\[
\mathrm{sq}_-(\mathscr C_\zeta)\ge 1.
\]
\end{corollary}
\begin{proof}
由定理~\ref{thm:xi-time-part9zq-threepoint-jensen-rh-disproof}，该见证推出 RH 失效。再由定理~\ref{thm:xi-disk-zero-pole-index-reduction}，
\[
\text{RH}\iff \mathrm{sq}_-(\mathscr C_\zeta)=0.
\]
故 RH 失效即推出
\[
\mathrm{sq}_-(\mathscr C_\zeta)\neq 0.
\]
又负平方指数取值于非负整数，遂有
\[
\mathrm{sq}_-(\mathscr C_\zeta)\ge 1.
\]
证毕。
\end{proof}

\begin{theorem}[主 primitive 扭曲的二次资源律]\label{thm:xi-time-part9zq-primitive-twist-quadratic-resource}
记
\[
\lambda:=\rho(1)=\varphi^2,
\qquad
r_p:=\frac{\rho(e^{2\pi i/p})}{\lambda}
\]
为奇素数模 \(p\) 上主 primitive 扭曲通道的归一化谱半径。定义
\[
A_\varphi:=\frac{6\sqrt5-5}{250}{(2\pi)}^2,
\qquad
B_\varphi:=\frac{7+24\sqrt5}{75000}{(2\pi)}^4.
\]
则当 \(p\to\infty\) 时，
\[
\log r_p=-A_\varphi p^{-2}+B_\varphi p^{-4}+O(p^{-6}),
\]
从而
\[
1-r_p
=
A_\varphi p^{-2}
\Bigl(-B_\varphi-\frac{A_\varphi^2}{2}\Bigr)p^{-4}
O(p^{-6}).
\]
特别地，若要求
\[
1-r_p\le \varepsilon,
\]
则必有
\[
p\ge \sqrt{\frac{A_\varphi}{\varepsilon}}\,(1+o(1))
\qquad
(\varepsilon\downarrow 0).
\]
因此，主 primitive 扭曲的抑制成本在素数模度量上必然服从二次资源律。
\end{theorem}
\begin{proof}
由推论~\ref{cor:real-input-40-collision-twist-fourth}，对
\[
t_p:=\frac{2\pi}{p}
\]
有
\[
\log\frac{\rho(t_p)}{\lambda}
=
-\frac{6\sqrt5-5}{250}\,t_p^2
+\frac{7+24\sqrt5}{75000}\,t_p^4
+O(t_p^6),
\]
即
\[
\log r_p=-A_\varphi p^{-2}+B_\varphi p^{-4}+O(p^{-6}).
\]
对其指数化，得到
\[
r_p
=
\exp\!\bigl(-A_\varphi p^{-2}+B_\varphi p^{-4}+O(p^{-6})\bigr)
=
1-A_\varphi p^{-2}
+\Bigl(B_\varphi+\frac{A_\varphi^2}{2}\Bigr)p^{-4}
+O(p^{-6}),
\]
从而
\[
1-r_p
=
A_\varphi p^{-2}
+\Bigl(-B_\varphi-\frac{A_\varphi^2}{2}\Bigr)p^{-4}
+O(p^{-6}).
\]
最后由
\[
1-r_p=A_\varphi p^{-2}+O(p^{-4})
\]
反解即可得到
\[
p\ge \sqrt{\frac{A_\varphi}{\varepsilon}}\,(1+o(1)).
\]
证毕。
\end{proof}

\begin{corollary}[四阶补偿把主 primitive 审计复杂度指数从 \texorpdfstring{$1/4$}{1/4} 降到 \texorpdfstring{$1/6$}{1/6}]\label{cor:xi-time-part9zq-primitive-twist-fourth-order-compensation}
沿用定理~\ref{thm:xi-time-part9zq-primitive-twist-quadratic-resource} 的记号。定义四阶补偿预测器
\[
\widehat r_p^{(4)}:=\exp\!\bigl(-A_\varphi p^{-2}+B_\varphi p^{-4}\bigr),
\]
以及二阶高斯预测器
\[
\widehat r_p^{(2)}:=\exp(-A_\varphi p^{-2}).
\]
则有
\[
\left|r_p-\widehat r_p^{(4)}\right|=O(p^{-6}),
\qquad
\left|r_p-\widehat r_p^{(2)}\right|=O(p^{-4}).
\]
因此，为达到误差阈值 \(\eta\)，四阶补偿仅需
\[
p\gg \eta^{-1/6},
\]
而二阶高斯近似则需
\[
p\gg \eta^{-1/4}.
\]
\end{corollary}
\begin{proof}
由定理~\ref{thm:xi-time-part9zq-primitive-twist-quadratic-resource}，
\[
r_p
=
\exp\!\bigl(-A_\varphi p^{-2}+B_\varphi p^{-4}+O(p^{-6})\bigr).
\]
故
\[
r_p-\widehat r_p^{(4)}
=
\widehat r_p^{(4)}\bigl(e^{O(p^{-6})}-1\bigr)
=
O(p^{-6}),
\]
而
\[
r_p-\widehat r_p^{(2)}
=
\widehat r_p^{(2)}\bigl(e^{B_\varphi p^{-4}+O(p^{-6})}-1\bigr)
=
O(p^{-4}).
\]
对这两个余项分别反解即得
\[
p\gg \eta^{-1/6},
\qquad
p\gg \eta^{-1/4}.
\]
证毕。
\end{proof}

\begin{theorem}[有限层 prime-slice 账本的活跃片数线性下界]\label{thm:xi-time-part9zq-active-primeslice-linear-growth}
设
\[
X_0\xrightarrow{\pi_0}X_1\xrightarrow{\pi_1}\cdots\xrightarrow{\pi_{k-1}}X_k
\]
为一条有限纤维满射塔，并假定每层都存在某个 \(y_j\in X_{j+1}\) 使
\[
\#\pi_j^{-1}(y_j)\ge 2
\qquad
(0\le j\le k-1).
\]
设
\[
\Psi(x_0)=\bigl(x_k,H(x_0)\bigr),
\qquad
H(x_0)=\prod_{j=0}^{k-1}h_j(x_0),
\]
其中每个 \(h_j(x_0)\) 的全部素因子均落在预先固定的 prime-slice \(\mathcal P_j\) 中，并设 \(\Psi\) 为单射。定义活跃片集合
\[
\mathcal A_k
:=
\left\{
0\le j\le k-1:\ \exists\,x_0\in X_0,\ h_j(x_0)\neq 1
\right\}.
\]
则
\[
|\mathcal A_k|\ge k.
\]
更强地，分层 prime-slice 最优构造在最坏情形下达到等号。
\end{theorem}
\begin{proof}
由假设，每一层都是非平凡分支层，故
\[
J=
\left\{
0,\dots,k-1
\right\}
\]
在定理~\ref{thm:xi-prime-slice-nontrivial-layer-exact-minimality} 的记号下满足 \(|J|=k\)。该定理给出任何这类按层分解的单射 prime-slice 外置都必须满足
\[
N_{\mathrm{slice}}(\Psi)\ge |J|=k.
\]
而这里
\[
N_{\mathrm{slice}}(\Psi)=|\mathcal A_k|,
\]
故
\[
|\mathcal A_k|\ge k.
\]

另一方面，定理~\ref{thm:xi-time-part53ac-primeslice-truncated-history-inverse-limit} 已给出恰在每层写入一枚有效素数的 sharp 构造，故在最坏情形下等号可达。证毕。
\end{proof}

\begin{corollary}[自然扩张一旦处处分支，活跃 prime-slice 必为无穷]\label{cor:xi-time-part9zq-natural-extension-infinite-primeslice-support}
在定理~\ref{thm:xi-time-part9zq-active-primeslice-linear-growth} 的设定下，设
\[
X_0\xleftarrow{\pi_0}X_1\xleftarrow{\pi_1}X_2\xleftarrow{\pi_2}\cdots
\]
是一条在每一深度都存在非平凡分支的自然扩张链。则任一与有限截断相容的 injective prime-slice 账本，在无限深度上必然具有无限活跃片支撑。特别地，若各 prime-slice \(\mathcal P_j\) 两两不交，则所用 prime-slice 的总数必为无穷。
\end{corollary}
\begin{proof}
对任意 \(k\ge 1\)，把账本限制到前 \(k\) 层截断，便得到定理~\ref{thm:xi-time-part9zq-active-primeslice-linear-growth} 的有限深度情形，因此前 \(k\) 层的活跃片数至少为 \(k\)。令 \(k\to\infty\)，便知无限深度上活跃片支撑不可能有限。

这也可视为定理~\ref{thm:xi-time-part63b-infinite-depth-finite-prime-register-impossibility} 的直接重述：一旦每层都真实分支，就不存在只用有限多个 prime-slice 的全历史精确外置。故若各 \(\mathcal P_j\) 两两不交，则所使用的 prime-slice 总数必为无穷。证毕。
\end{proof}

\begin{theorem}[固定有限支持局部化 solenoid 不可能忠实承载无穷分支历史]\label{thm:xi-time-part9zq-no-fixed-finite-support-solenoid-host}
设 \(S\subset\mathbb P\) 为有限素数集，并记
\[
G_S:=\ZZ[S^{-1}],
\qquad
\Sigma_S:=\widehat{G_S}.
\]
考虑一条在每一深度都存在非平凡分支的自然扩张链。则不存在任何忠实的 prime-support 外置模型，能够把其全部无穷分支历史压入某个固定有限支持的局部化系统 \(G_S\) 或其 Pontryagin 对偶 \(\Sigma_S\) 中。
\end{theorem}
\begin{proof}
由定理~\ref{thm:killo-localization-circle-dimension-prime-ledger-splitting}，固定有限支持 \(S\) 的 Pontryagin 对偶满足短正合列
\[
0\to \prod_{p\in S}\ZZ_p\to \widehat{G_S}\to \TT\to 0,
\]
其全不连通核的 prime-support 恰为有限集 \(S\)，并且 \(S\) 可由该核完全恢复。故任何通过 \(G_S\) 或 \(\Sigma_S\) 实现的 prime-support 外置模型，至多只能动用有限多个 prime-slice 方向。

另一方面，推论~\ref{cor:xi-time-part9zq-natural-extension-infinite-primeslice-support} 表明：只要自然扩张在每一深度都有真实分支，则任何忠实账本都必须具有无限活跃片支撑。二者矛盾，故固定有限支持的局部化 solenoid 不可能忠实承载全部无穷分支历史。证毕。
\end{proof}

\noindent
由此，Jensen 递归阈值、primitive twist 慢模与 prime-slice/solenoid 外置便在同一条有限层刚性链上闭合：中心单零点在 dyadic 细化中存在 sharp 的触发与自证阈值；主 primitive twist 的抑制代价被 \(p^{-2}\) 资源律刚性锁定；而一切处处分支的前史链又不能被任何固定有限支持的 prime-ledger 或 solenoid 宿主忠实容纳。若继续沿 RH 路线推进，则更自然的工作方向不再是寻找单一有限宿主去同时吸收这些现象，而是分别建立 Jensen 阈值、twist 资源律与 prime-support 外置障碍之间的统一结构理论。

\endinput
